% M04b_dft_numerics.tex -- part 2 of ch:mono-dft-theory: basis, PAW and band counting, SCF and force errors, forces
% and relaxation, slabs and dipoles, band alignment, computational cost. Numerical checks use raw outputs:
% plane-wave counts and RAM estimates from cern_htcondor/results/{slab2_relax_c2,slab1r_final_cpu,iwo_W24d}/scf*.out and
% bulk_In2O3_protocol/{vcrelax_ecut71_k3,bands_path}.out (read 2026-10-14); force-error figures from ledger C034/C035.

% =====================================================================================================================
\section{Plane waves and the cutoff}
\label{sec:mono-dft-pw}

Because $u_{n\vk}$ is lattice-periodic it is a Fourier series over reciprocal-lattice vectors $\vG$,
$\psi_{n\vk}(\vr) = \Omega^{-1/2}\sum_{\vG} c_{n,\vk+\vG}\,e^{i(\vk+\vG)\cdot\vr}$. The series is truncated at a kinetic
cutoff, $\tfrac12|\vk+\vG|^2 \le \Ecut$ (in Rydberg units, $|\vk+\vG|^2 \le \Ecut$). The number of plane waves is the
number of reciprocal-lattice points in a sphere of radius $k_c = \Ecut^{1/2}$ (bohr$^{-1}$, $\Ecut$ in Ry), each
occupying a volume $(2\pi)^3/\Omega$:
\begin{equation}
N_{\mathrm{pw}} \approx \frac{4\pi k_c^3/3}{(2\pi)^3/\Omega} = \frac{\Omega\,\Ecut^{3/2}}{6\pi^2} .
\label{eq:mono-dft-npw}
\end{equation}
\Cref{tab:mono-dft-npw} checks \cref{eq:mono-dft-npw} against the plane-wave counts printed by \texttt{pw.x} at
$\Gamma$ for $\Ecut = 71$~Ry: the agreement is better than 0.15\,\%, so the cost model below rests on the actual basis
sizes. The density and potential need components up to $2k_c$, i.e.\ a density cutoff of $4\Ecut$ for norm-conserving
pseudopotentials; PAW augmentation charges are harder, and the project uses $8\Ecut$ (568~Ry), which multiplies the
number of density components by $8^{3/2} \approx 22.6$.

\begin{table}[htbp]\centering\small
\caption{Plane-wave basis at 71~Ry: \cref{eq:mono-dft-npw} against the count printed at $\Gamma$ (raw outputs).}
\label{tab:mono-dft-npw}
\begin{tabular}{lrrrr}\toprule
Cell & $\Omega$ (bohr$^3$) & \cref{eq:mono-dft-npw} & printed & bands \\\midrule
bulk \InO{}, 40-atom primitive (vc-relaxed) & 3\,693.6 & 37\,315 & 37\,359 & -- \\
bulk In$_{31}$WO$_{48}$, 80 atoms & 7\,387.2 & 74\,630 & 74\,653 & 424 \\
1~nm slab In$_{24}$O$_{48}$H$_{24}$, 96 atoms & 25\,140.0 & 253\,982 & 253\,973 & 328 \\
2~nm slab In$_{56}$O$_{96}$H$_{24}$, 176 atoms & 32\,429.2 & 327\,622 & 327\,585 & 680 \\\bottomrule
\end{tabular}
\end{table}

The cutoff of 71~Ry met the pre-registered criteria for the primitive bulk cell (\cref{sec:dft-calc-conv}); the
amorphous relaxations use 60~Ry for cost and the amorphous final SCFs return to 71~Ry, so electronic read-outs share the
crystalline basis quality. Because plane waves do not depend on atomic positions, the Hellmann--Feynman forces need no
Pulay correction; the stress does, through the change of $N_{\mathrm{pw}}$ with the cell, which is why variable-cell
relaxations use a modified kinetic functional or a higher cutoff.

% =====================================================================================================================
\section{Pseudopotentials, PAW and the counting of bands}
\label{sec:mono-dft-paw}

Core electrons are frozen. The projector augmented-wave method \citep{Bloechl1994,Kresse1999} maps smooth
pseudo-orbitals $|\tilde\psi\rangle$ onto all-electron orbitals with a linear transformation,
\begin{equation}
|\psi\rangle = |\tilde\psi\rangle + \sum_i \big(|\phi_i\rangle - |\tilde\phi_i\rangle\big)\langle\tilde p_i|\tilde\psi\rangle ,
\label{eq:mono-dft-paw}
\end{equation}
where $\phi_i$ and $\tilde\phi_i$ are all-electron and pseudo partial waves inside spheres around the atoms and
$\tilde p_i$ dual projectors. The data sets of this project (pslibrary 1.0.0, \citealp{DalCorso2014}) put the In 4d
shell and the W 5s5p shells into the valence. This choice has a consequence that caused the most important correction
of the project \corr{X1}: the number of bands below the gap is not the same in a pure and a W-doped cell.

\paragraph{Band counting.} With valence charges $Z_{\mathrm{In}} = 13$ (4d$^{10}$5s$^2$5p$^1$), $Z_{\mathrm{W}} = 14$
(5s$^2$5p$^6$5d$^4$6s$^2$), $Z_{\mathrm{O}} = 6$ and $Z_{\mathrm{H}} = 1$, the occupied manifold below the gap of an
In--W--O--H cell consists of five In 4d bands per In, four W semicore bands (5s, 5p) per W, and four O 2s/2p bands per
O (the O--H bonds lie inside the O manifold). Hence
\begin{equation}
N_{\mathrm{VB}} = 5n_{\mathrm{In}} + 4n_{\mathrm{W}} + 4n_{\mathrm{O}}, \qquad
N_{\mathrm{el}} = 13n_{\mathrm{In}} + 14n_{\mathrm{W}} + 6n_{\mathrm{O}} + n_{\mathrm{H}},\qquad
N_{\mathrm{ex}} = N_{\mathrm{el}} - 2N_{\mathrm{VB}} .
\label{eq:mono-dft-bands}
\end{equation}
For the cells of this project: In$_{32}$O$_{48}$: $N_{\mathrm{el}} = 704$, $N_{\mathrm{VB}} = 352$, $N_{\mathrm{ex}} = 0$;
In$_{31}$WO$_{48}$: $705$, $351$, $3$; In$_{24}$O$_{48}$H$_{24}$: $624$, $312$, $0$; In$_{23}$WO$_{48}$H$_{24}$: $625$,
$311$, $3$. The electron counts are those printed by \texttt{pw.x}. Substituting one W for one In removes one band from
the occupied manifold and adds three electrons above it: the valence-band top of a W cell is band $n_{\mathrm{pure}}-1$,
and the three electrons that W donates occupy bands $n_{\mathrm{pure}}$ and above. Reading band $n_{\mathrm{pure}}$ as the
valence-band top, as the analysis of 30 September did, places a conduction-band state at the valence-band top and
produces exactly the spurious ``W level below the conduction band'' that was corrected on 5 October. The amorphous
read-out avoids the problem by construction: it defines $N_v = (N_{\mathrm{el}} - N_{\mathrm{ex}})/2$ from the nominal
excess count and checks the occupation above $N_v$ (\cref{sec:mono-am-defs}).

% =====================================================================================================================
\section{The self-consistent field and the size of force errors}
\label{sec:mono-dft-scf}

\Cref{eq:mono-dft-KS} is a fixed-point problem $n = \mathcal{F}[n]$, solved by density mixing. Linear mixing,
$n_{i+1} = n_i + \beta(\mathcal{F}[n_i]-n_i)$, converges only if $\beta$ is smaller than about the inverse of the
largest eigenvalue of the dielectric response; in metallic or large cells long-wavelength charge oscillations
(``charge sloshing'') make that eigenvalue scale with the square of the cell length. Broyden-type mixing
\citep{Johnson1988} with Kerker \citep{Kerker1981} or local Thomas--Fermi preconditioning damps these modes. This is why
the metallic 2~nm films needed $\beta = 0.1$ and \texttt{conv\_thr} $10^{-8}$, and why the amorphous rescues
(\amend{A15}) added smaller $\beta$, more mixing vectors and local-TF mixing.

\paragraph{Energy and force errors.} The Kohn--Sham energy is variational in the density: an SCF error $\delta n$
changes $E$ only to second order, $\delta E \propto \|\delta n\|^2$. The Hellmann--Feynman force is not variational, so
its error is first order, $\delta F \propto \|\delta n\|$. Quantum ESPRESSO's \texttt{conv\_thr} bounds the estimated
energy error, so the force error scales as
\begin{equation}
\delta F \propto \sqrt{\texttt{conv\_thr}} .
\label{eq:mono-dft-forceerr}
\end{equation}
The logged amorphous checks test this. In two 1.0~nm amorphous cells of similar size (134 and 138 atoms) the printed
``total SCF correction'' to the forces is $1.3\times10^{-3}$~Ry/bohr at the production threshold $10^{-7}$
(\claim{C034}) and $3.3$--$3.7\times10^{-7}$~Ry/bohr at the check threshold $10^{-14}$ (\claim{C035}). The ratio, about
3700, is close to $\sqrt{10^{7}} \approx 3200$, as \cref{eq:mono-dft-forceerr} predicts. The production force noise,
about $10^{-3}$~Ry/bohr, is of the same size as the relaxation threshold itself. This is why the project compares
fixed-geometry forces at tightened thresholds when it needs force differences smaller than $10^{-3}$, and why a
restart can fail to reproduce a continued SCF at production tolerance (\claim{C036}).

\section{Forces and relaxation}
\label{sec:mono-dft-forces}

For an exact eigenstate, $\partial E/\partial\vR_I = \langle\Psi|\partial\Hop/\partial\vR_I|\Psi\rangle$ because the
terms containing $\partial\Psi/\partial\vR_I$ vanish by normalisation (Hellmann--Feynman theorem, \citealp{Feynman1939}).
In DFT with a position-independent basis this gives $\vec F_I = -\int n(\vr)\,\partial\vext/\partial\vR_I\,d\vr -
\partial E_{II}/\partial\vR_I$, plus PAW and smearing terms that \texttt{pw.x} includes. Relaxations use BFGS
\citep{Billeter2003} and stop when every force component is below \texttt{forc\_conv\_thr} ($10^{-3}$~Ry/bohr) and the
energy change below \texttt{etot\_conv\_thr} ($10^{-4}$~Ry). The criterion is per component: the relaxed pure 1~nm slab
satisfies it while its total force, the norm over all 288 components, is $6.1\times10^{-3}$~Ry/bohr, which corresponds
to a typical component of $6.1\times10^{-3}/\sqrt{288} \approx 3.6\times10^{-4}$~Ry/bohr. Quantum ESPRESSO prints final
coordinates also after a failed BFGS run, which is why the amorphous chains count a relaxation as converged only after
the ``bfgs converged'' line (\amend{A19}). Variable-cell relaxations (Wentzcovitch dynamics, \citealp{Wentzcovitch1991})
relax the bulk lattice and, for slabs, only the in-plane cell (\texttt{cell\_dofree = '2Dxy'}).

% =====================================================================================================================
\section{Slabs: vacuum, planar averages and dipoles}
\label{sec:mono-dft-slab}

A film is modelled as a slab repeated along $z$ with vacuum in between. The planar average of the electrostatic
potential, $\bar V(z) = A^{-1}\int_A V(x,y,z)\,dx\,dy$, approaches a constant in the vacuum, the vacuum level
$\Evac$; the electron affinity and ionisation potential follow as $\mathrm{EA} = \Evac - E_{\mathrm{CBM}}$ and
$\mathrm{IP} = \Evac - E_{\mathrm{VBM}}$ (\cref{eq:dft-eaip}). Inside the slab $\bar V$ oscillates with the atomic layers;
the macroscopic average $\bar{\bar V}(z) = l^{-1}\int_{z-l/2}^{z+l/2}\bar V(z')\,dz'$ over one period $l$ removes the
oscillation and leaves the slowly varying part \citep{Baldereschi1988}.

\paragraph{Polar slabs and the dipole correction.} If the two faces of a slab differ, the slab carries a dipole moment
per area $p$. Its potential jumps by $\Delta\phi = p/\varepsilon_0$ across the slab. Periodicity along $z$ forces the
average field over the cell to vanish, so the calculation adds a uniform compensating field $\mathcal{E} =
\Delta\phi/c$ over the whole cell of height $c$; in the vacuum this appears as a linear ramp of $\bar V$ instead of a
plateau, and the field also acts on the slab. A dipole correction \citep{Neugebauer1992,Bengtsson1999} adds a sawtooth
potential that cancels the field and places the discontinuity in the vacuum. The amorphous slabs use it; the
crystalline slabs do not, because the 1 and 2~nm films contain an operation that maps one face onto the other.

The 0.695~nm crystalline film is polar (\cref{sec:mono-dft-thin}). Inverting the relation above with its logged ramp
$-13.6$~meV/\AA{} and its cell height $c = 32.48$~\AA{} gives a potential step $\Delta\phi = 0.44$~V and a dipole of
about $0.26\,e$\,\AA{} per $107$~\AA$^2$ cell ($1.3$~D) \fact{019}. The compensating field drops about 0.095~eV across
the 0.695~nm film. That is the scale of the (untested) effect of the artificial field on the gap: of order 0.1~eV
against a gap of 2.2~eV, and of the same order as the reference-choice sensitivity of $\pm0.17$~eV on $\dEc$
\fact{018}. A dipole-corrected SCF would replace this estimate.

\section{Band alignment between slab and bulk}
\label{sec:mono-dft-align}

A bulk calculation has no vacuum, so its band edges are known only relative to its own average potential. The two-step
method \citep{VandeWalle1987,Baldereschi1988} combines (i) the bulk band edges relative to the bulk macroscopic average,
$E_{c}^{\mathrm{bulk}} - \bar{\bar V}^{\mathrm{bulk}}$, with (ii) the slab interior macroscopic average relative to the
vacuum, $\bar{\bar V}^{\mathrm{slab}}_{\mathrm{int}} - \Evac$, to place bulk-like edges on the vacuum scale:
\begin{equation}
\dEc(t) = \big[E_{\mathrm{CBM}}^{\mathrm{slab}} - \Evac\big] - \big[E_{c}^{\mathrm{bulk}} - \bar{\bar V}^{\mathrm{bulk}}\big]
 - \big[\bar{\bar V}^{\mathrm{slab}}_{\mathrm{int}} - \Evac\big],
\label{eq:mono-dft-twostep}
\end{equation}
and likewise for $\Delta E_{\mathrm{v}}$ (signed; negative means downwards), with $\Delta\Eg = \dEc - \Delta E_{\mathrm{v}}$.
The method needs an interior region in which $\bar{\bar V}$ is flat. In the 2~nm film the averaging window of half the
lattice constant, the period of bixbyite along [001], gives $\dEc = +0.281$~eV and $\Delta E_{\mathrm{v}} = -0.054$~eV
(share of the opening in the conduction band 0.84) \fact{014}. A 1~nm film has no such interior; its values follow from
the 2~nm result plus the differences of EA and IP between the two films, which share the same termination, giving
$+0.848$ and $-0.052$~eV (share 0.94) \fact{015}. The amorphous read-out applies the same rule: alignment is valid only
with an interior plateau with a standard deviation of at most 0.05~eV (\cref{sec:mono-am-defs}).

% =====================================================================================================================
\section{Computational cost}
\label{sec:mono-dft-cost}

For a cell of $N$ atoms, $N_{\mathrm{pw}} \propto \Omega \propto N$ and the number of bands $N_b \propto N$. The
dominant operations of an iterative plane-wave solver are the application of $\Hop$ by fast Fourier transforms,
$O(N_bN_{\mathrm{pw}}\log N_{\mathrm{pw}})$, and orthonormalisation and subspace rotation, $O(N_b^2N_{\mathrm{pw}})$; the
latter dominates for large cells, so the cost per SCF iteration grows as $N^3$ and the memory for the wave functions as
$N_bN_{\mathrm{pw}}N_k \propto N^2N_k$. Linear-scaling methods exist \citep{Goedecker1999} but were not used. The
numbers of this project illustrate the scaling. Going from the 1~nm to the 2~nm slab doubles the bands (328 to 680) and
increases the plane waves by 29\,\% (\cref{tab:mono-dft-npw}): the wave functions grow by a factor of 2.7 and the
$N_b^2N_{\mathrm{pw}}$ work by 5.5. One set of wave functions of the 2~nm slab at one $k$ point occupies
$680\times327\,585\times16$~bytes $=3.6$~GB, and the four irreducible $k$ points with the Davidson workspace and the
density grids bring the run to the 82~GB per process that \texttt{pw.x} estimates for the single-process GPU relaxation;
this is why the 176-atom film needed a large-memory GPU or a distributed CPU run, and why the 250--300-atom amorphous
cells fitted on one 96~GB GPU only with split species labels and reduced cutoffs during relaxation
(\cref{sec:mono-am-relax}). A hybrid functional adds an exact-exchange term that couples all pairs of occupied orbitals
at $k$ and $k+q$, multiplying the cost by roughly $N_qN_b$ FFT pairs per orbital; this is why hybrid calculations were
limited to fixed-geometry checks of the bulk and the 1~nm film on CPUs (\cref{sec:dft-hse}).

\begin{keyresult}
The first-principles numbers of this monograph rest on five approximations whose size can be stated: the adiabatic
separation (negligible for the quantities used); the PBE functional (gaps too small by a factor of about three, which is
why only confinement shifts are transferred and checked with HSE06, ratio 1.06--1.11); finite bases and $k$ grids
(converged by registered criteria; the plane-wave counts reproduce the analytic estimate within 0.15\,\%); smearing
(physical occupations only for Gaussian smearing well below the energy scales of interest; Marzari--Vanderbilt weights
can be negative); and SCF tolerance (force noise of about $10^{-3}$~Ry/bohr at production settings, scaling as the
square root of the energy threshold, confirmed by the logged checks). The counting of semicore bands explains the
project's main correction \corr{X1} exactly.
\end{keyresult}
